% =============================================================================
% CHAPTER 2 — THEORETICAL FOUNDATION (REDUCED, Level B)
% Thesis: PATHCAST
% Target: ~5,000 words / ~18-22 pages
% Strategy: keep formal definitions and main results; defer derivations,
%           proofs, algorithmic details, and numerical examples to appendices
%           G (Markov), H (Probabilistic Metrics), I (Process Mining details),
%           J (End-to-end walkthrough).
% =============================================================================

\chapter{Theoretical Foundation}
\label{ch:theoretical-foundation}

This chapter presents the theoretical pillars on which PATHCAST is built:
process mining (Section~\ref{sec:process-mining}), mining of software
repositories (Section~\ref{sec:msr}), discrete-time Markov chains
(Section~\ref{sec:markov}), Monte Carlo simulation
(Section~\ref{sec:monte-carlo}), probabilistic forecasting metrics
(Section~\ref{sec:prob-metrics}), and machine learning for software
engineering (Section~\ref{sec:ml-se}). Section~\ref{sec:related-work}
synthesizes the gap addressed by this thesis. Section~\ref{sec:ch2-synthesis}
integrates the five domains into a unified conceptual framework. Extensive
derivations, algorithmic details, and worked examples are deferred to
Appendices~G--J.


% =============================================================================
% 2.1 PROCESS MINING
% =============================================================================
\section{Process Mining}
\label{sec:process-mining}

Process mining extracts knowledge from event logs recorded by information
systems~\cite{aalst2016process}. Unlike interview-based modeling, it
reconstructs actual process behavior from execution data, bridging data
science and process science.

\subsection{Event Logs}
\label{sec:event-logs}

\begin{definition}[Event]
\label{def:event}
An event $e = (c, a, t, d)$ has a case identifier $c \in \mathcal{C}$, an
activity name $a \in \mathcal{A}$, a timestamp $t \in \mathbb{R}^+$, and an
attribute map $d \in \mathcal{D}$. Projection functions $\kappa(e)$,
$\alpha(e)$, $\tau(e)$ extract the case, activity, and timestamp.
\end{definition}

\begin{definition}[Trace and Event Log]
\label{def:event-log}
A trace $\sigma_c = \langle e_1, \ldots, e_n \rangle$ is the timestamp-ordered
event sequence of case $c$. An event log
$\mathcal{L} = \{\sigma_{c_1}, \ldots, \sigma_{c_N}\}$ is a multiset of
traces. The activity set is
$\mathcal{A}_\mathcal{L} = \bigcup_{e \in \mathcal{L}} \{\alpha(e)\}$.
\end{definition}

The IEEE 1849-2016 XES standard~\cite{xes2016} formalizes the log structure
through the levels \emph{log}, \emph{trace}, and \emph{event}, with
extensions \texttt{concept:name}, \texttt{time:timestamp},
\texttt{org:resource}, and \texttt{lifecycle:transition}. Object-Centric
Event Logs (OCEL)~\cite{vanderAalst2019ObjectCentric} generalize the
single-case model to many-to-many object relationships, which is relevant
for software development where one commit may relate to multiple issues.

\subsection{The Three Types of Process Mining}
\label{sec:pm-types}

Following van der Aalst~\cite{aalst2016process}: \emph{Discovery} extracts a
model $M$ from $\mathcal{L}$ alone; \emph{Conformance Checking} compares
$\mathcal{L}$ against a given $M$; \emph{Enhancement} enriches $M$ with
information derived from $\mathcal{L}$ (e.g., performance overlays). PATHCAST
uses all three: discovery to extract the state-transition backbone,
conformance to validate the backbone and produce per-case features, and
enhancement to obtain sojourn time distributions.

\subsection{Process Discovery and Quality Criteria}
\label{sec:process-discovery}

Discovery algorithms transform an event log into a process model. The main
families relevant for software event logs are summarized in
Table~\ref{tab:discovery-algorithms}. Algorithmic internals and complexity
proofs are presented in Appendix~I.

\begin{table}[htbp]
\centering
\caption{Process discovery algorithms used in software process mining.}
\label{tab:discovery-algorithms}
\small
\begin{tabular}{lcccc}
\toprule
\textbf{Criterion} & \textbf{Alpha} & \textbf{Heuristics} &
\textbf{Inductive} & \textbf{Split} \\
\midrule
Output formalism     & Petri net & C-net     & Process tree & BPMN \\
Noise handling       & None      & Threshold & Filtering    & Filtering \\
Short loops          & No        & Yes       & Yes          & Yes \\
Soundness guarantee  & No        & No        & Yes          & No \\
\bottomrule
\end{tabular}
\end{table}

PATHCAST adopts the \textbf{Inductive Miner} (IM)~\cite{leemans2013discovering}
because it is the only algorithm in the table that guarantees
\emph{soundness} (no deadlocks) by construction, producing a block-structured
process tree that maps naturally onto the hierarchical structure of SDLC
workflows. The IMf variant filters infrequent behavior through a noise
threshold $\theta_{\text{noise}} \in [0,1]$.

The quality of a discovered model is evaluated along four
dimensions~\cite{aalst2016process}: \emph{fitness} (capacity to replay the
log), \emph{precision} (no behavior allowed beyond what was observed),
\emph{generalization} (capacity to represent unseen but legitimate behavior),
and \emph{simplicity} (parsimony). The first three are formalized below; the
fourth is constrained by the SDLC taxonomy and not reported separately.

\subsection{Conformance Checking}
\label{sec:conformance-checking}

\begin{definition}[Alignment]
\label{def:alignment}
An alignment $\gamma$ between trace $\sigma$ and model $M$ is a sequence of
moves, each classified as \emph{synchronous} ($s_i = m_i$),
\emph{log move} ($m_i = \gg$), or \emph{model move} ($s_i = \gg$). The
\emph{optimal alignment} $\gamma^*$ minimizes the total non-synchronous
cost~\cite{adriansyah2011alignments}.
\end{definition}

The optimal alignment is computed via $A^*$ on the synchronous product of
trace and model reachability graphs; full algorithmic details are in
Appendix~I.

\begin{definition}[Alignment-Based Fitness]
\label{def:fitness}
\begin{equation}
\label{eq:fitness}
\text{fitness}(\sigma, M) = 1 - \frac{\text{cost}(\gamma^*)}
{\text{cost}(\gamma_{\text{worst}})}
\end{equation}
Log-level fitness is the trace-weighted average.
\end{definition}

\begin{definition}[Case-Level Conformance]
\label{def:case-conformance}
\begin{equation}
\label{eq:case-conformance-cap2}
\text{conf}(c) = \frac{|\{(s_i, m_i) \in \gamma^*_c : s_i = m_i\}|}{|\gamma^*_c|}
\end{equation}
\end{definition}

Case-level conformance is used in Chapter~\ref{ch:method} as a feature for
the ML component. \emph{Precision} is computed via the ETC (Escaping Edges)
metric and \emph{generalization} via state-visit frequency analysis; both
complement fitness for full quality assessment.

\subsection{Process Enhancement}
\label{sec:process-enhancement}

Enhancement annotates the discovered model with temporal information
extracted from the log: \emph{sojourn time} (total time per state visit),
\emph{waiting time} (time before processing starts), \emph{service time}
(active processing), and \emph{throughput time} (end-to-end case duration).
For PATHCAST, enhancement produces the empirical sojourn distributions
$\hat{F}_s$ that parameterize the Monte Carlo simulation
(Section~\ref{sec:monte-carlo}).

\subsection{Predictive Process Monitoring}
\label{sec:ppm}

Predictive Process Monitoring (PPM)~\cite{tax2017,marquezChamorro2018}
predicts the next activity, remaining time, or final outcome of a running
case from its partial trace. PPM is conceptually adjacent to PATHCAST but
differs in scope: PPM operates at the individual case level, typically
treating the process as a black box (especially with deep-learning
methods); PATHCAST operates at the process level, generating distributional
forecasts via explicit state-transition modeling. The two are complementary
rather than competing.


% =============================================================================
% 2.2 MINING SOFTWARE REPOSITORIES
% =============================================================================
\section{Mining Software Repositories}
\label{sec:msr}

Mining Software Repositories (MSR) extracts knowledge from software
development artifacts: version control, issue trackers, code review
platforms, CI/CD pipelines, and communication archives~\cite{hassan2008road,
kagdi2007}. The bridge between MSR and process mining was established by
Cook and Wolf~\cite{cook1998} and formalized by Poncin et
al.~\cite{poncin2011process}: instead of analyzing artifacts in isolation,
software events are reframed as sequences belonging to process instances.

\subsection{Data Sources}
\label{sec:msr-data-sources}

Modern software development emits events across heterogeneous, interconnected
tools:
\begin{itemize}
  \item \textbf{Issue trackers} (Jira, GitHub Issues): status transitions
  with timestamp, actor, and contextual attributes.
  \item \textbf{Version control} (Git, SVN): commit hash, author,
  timestamp, modified files.
  \item \textbf{Code review} (GitHub PRs, Gerrit): PR creation, reviews,
  approvals, merges.
  \item \textbf{CI/CD} (Jenkins, GitHub Actions): build, test, and deploy
  outcomes with duration and triggering commit.
\end{itemize}

The challenge is integration: a unified event log requires identifier
mapping, temporal alignment, and semantic harmonization across
tools~\cite{kochhar2019moving, Dixit2021}. The case correlation problem
is the entry point of PATHCAST Stage~1 (Chapter~\ref{ch:method}).

\subsection{Repository Dataset Infrastructure}

Empirical studies rely on curated datasets and collection tooling: SEART
GitHub Search~\cite{seart2021}, GH~Archive~\cite{gharchive2023},
RepoReapers~\cite{munaiah2017reporeapers},
SmartSHARK~\cite{trautsch2020smartshark},
GrimoireLab~\cite{dueenas2021grimoirelab}, and the Public Jira
Dataset~\cite{montgomery2022jira}. These reduce candidate-generation cost
but do not eliminate the validity threats documented by
Kalliamvakou~\cite{kochhar2019moving}: open-source bias, English-language
bias, GitHub-hosted bias, bot-generated noise, and incompleteness of
observable activity.

\subsection{Event Log Construction}

Construction involves three sub-tasks:
\begin{enumerate}
  \item \textbf{Event extraction}: tool-specific adapters retrieving
  timestamped records.
  \item \textbf{Case correlation}: mapping events to cases through explicit
  references (issue keys), structural links (PR-issue, build-commit), or
  heuristics (temporal proximity, author).
  \item \textbf{Activity mapping}: $\mu: \mathcal{A}_{\text{raw}} \to
  \mathcal{A}_{\text{standard}}$ projecting source-specific labels onto
  a canonical SDLC vocabulary.
\end{enumerate}

\subsection{Software Process Archaeology}

\emph{Software archaeology}~\cite{Zimmermann2005MiningVersionHistories}
denotes the reconstruction of process knowledge from legacy artifacts.
PATHCAST treats archaeological reconstruction as an upstream evidential
prerequisite, not as an independent objective. Three archaeological
principles are operationally relevant: \emph{stratigraphy} (events from
distinct periods reflect distinct process configurations),
\emph{contextual interpretation} (events must be interpreted against
project conventions), and \emph{taphonomic bias} (unobservable activities
bound forecast completeness).


% =============================================================================
% 2.3 MARKOV CHAINS
% =============================================================================
\section{Stochastic Process Modeling with Markov Chains}
\label{sec:markov}

Software delivery systems evolve through discrete states (Backlog,
In Progress, Review, Testing, Done). Queueing, rework, blocking, and
variable service times make these systems stochastic. Discrete-time Markov
chains (DTMCs) provide the mathematical formalism~\cite{Spedicato2017,
Haggstrom2002}. Detailed proofs and derivations are in Appendix~G.

\subsection{Discrete-Time Markov Chains}
\label{sec:dtmc}

\begin{definition}[Discrete-Time Markov Chain]
\label{def:dtmc}
A stochastic process $\{X_t\}_{t \in \mathbb{N}}$ on finite state space
$S = \{1, \ldots, m\}$ is a DTMC if for all $t$ and $i_0, \ldots, i_t, j$:
\begin{equation}
\label{eq:markov-property}
P(X_{t+1} = j \mid X_t, \ldots, X_0) = P(X_{t+1} = j \mid X_t)
\end{equation}
A time-homogeneous DTMC is fully described by a single transition matrix.
\end{definition}

\begin{definition}[Transition Matrix]
\label{def:transition-matrix}
$P = [p_{ij}]$ with $p_{ij} = P(X_{t+1}=j \mid X_t=i)$ is row-stochastic:
$p_{ij} \geq 0$ and $\sum_j p_{ij} = 1$ for all $i$.
\end{definition}

The Chapman-Kolmogorov relation $P^{(n)} = P^n$ enables multi-step
transition computation from a single estimated matrix. State distributions
evolve as $\pi^{(t)} = \pi^{(0)} P^t$.

\begin{definition}[Stationary Distribution]
\label{def:stationary}
A row vector $\pi$ with $\pi = \pi P$. For finite, irreducible, aperiodic
DTMCs the stationary distribution is unique and represents the long-run
state-occupancy fractions.
\end{definition}

States are classified as \emph{transient} (positive escape probability) or
\emph{recurrent} (return with probability 1). A chain with
\emph{absorbing} states ($p_{aa}=1$) reachable from every transient state
is the case relevant for SDLC forecasting.

\subsection{Absorbing Chains and the Fundamental Matrix}
\label{sec:absorbing-chains}

\begin{definition}[Absorbing Chain Canonical Form]
\label{def:absorbing}
For an absorbing chain with $t$ transient and $r$ absorbing states:
\begin{equation}
\label{eq:absorbing-form}
P = \begin{pmatrix} Q & R \\ \mathbf{0} & I_r \end{pmatrix}
\end{equation}
\end{definition}

\begin{definition}[Fundamental Matrix]
\label{def:fundamental-matrix}
\begin{equation}
\label{eq:fundamental}
N = (I_t - Q)^{-1} = \sum_{k=0}^{\infty} Q^k
\end{equation}
$N_{ij}$ is the expected number of visits to transient state $j$
starting from transient state $i$, before absorption. The matrix is
well-defined because the spectral radius of $Q$ is strictly less than~1.
\end{definition}

\begin{theorem}[Properties of Absorbing Chains]
\label{thm:absorption}
Let $N=(I-Q)^{-1}$. Then:
\emph{(i)} $E[T_i] = (N\mathbf{1})_i$ is the expected steps to absorption;
\emph{(ii)} $\text{Var}(T_i) = (N(2N_{\text{dg}}-I)\mathbf{1})_i - E[T_i]^2$;
\emph{(iii)} $B = NR$ is the absorption probability matrix.
\end{theorem}

\noindent The proof and derivations of (i)--(iii), together with
invertibility and convergence properties, are in Appendix~G.

\subsection{Sojourn Time Distributions}
\label{sec:sojourn-time}

\begin{definition}[Sojourn Time]
\label{def:sojourn-time}
For consecutive events $e_k, e_{k+1}$ with $\alpha(e_k) \neq \alpha(e_{k+1})$,
the sojourn time in state $s = \alpha(e_k)$ is
$\Delta t_s = \tau(e_{k+1}) - \tau(e_k)$.
\end{definition}

The empirical distribution
$\hat{F}_s = \{\Delta t_s^{(1)}, \ldots, \Delta t_s^{(n_s)}\}$ aggregates
all observed sojourns for state $s$. Software development sojourn times
are typically right-skewed and heavy-tailed; PATHCAST samples directly
from $\hat{F}_s$ rather than fitting a parametric family, avoiding
misspecification (Remark~\ref{rem:nonparametric} in
Chapter~\ref{ch:method}). Mean sojourn time is
$\bar{T}_s = \tfrac{1}{n_s} \sum_{k=1}^{n_s} \Delta t_s^{(k)}$
(formal definition in Equation~\ref{eq:mean-sojourn},
Chapter~\ref{ch:method}).

The augmentation of a DTMC with state-dependent sojourn distributions
yields a \emph{semi-Markov process}~\cite{RoggeSolti2015NonMarkovianSPN};
the formal connection is established in Appendix~G.

\begin{definition}[Semi-Markov Process]
\label{def:semi-markov}
A semi-Markov process is a pair (DTMC $\{X_n\}$, sojourn distributions
$\{H_{ij}(\cdot)\}$) where the process remains in state $X_n$ for a random
time drawn from $H_{X_n, X_{n+1}}$ before transitioning.
\end{definition}

\subsection{Estimation from Event Logs}
\label{sec:markov-estimation}

Given counts $c_{ij}$ of observed $i \to j$ transitions, the maximum
likelihood estimator is:
\begin{equation}
\label{eq:mle-transition}
\hat{p}_{ij} = \frac{c_{ij}}{\sum_{k} c_{ik}}
\end{equation}

Laplace smoothing with pseudo-count $\alpha > 0$ addresses the zero-count
problem; under a symmetric Dirichlet prior, the posterior is conjugate.
Smoothing details and the Bayesian credible-interval derivation are in
Appendix~G.

\subsection{Limitations of the Markov Assumption}
\label{sec:markov-limitations}

The first-order Markov assumption ignores history dependence (rework
escalation, defect severity) and time-varying transitions (team learning,
release cycles). PATHCAST mitigates these limitations through two
mechanisms: (i)~optional state-space expansion to encode history (e.g.,
distinguishing first-time vs.\ reworked Testing), and (ii)~the ML component
(Section~\ref{sec:ml-se}), which captures non-Markovian features as residual
predictive signal.


% =============================================================================
% 2.4 MONTE CARLO SIMULATION
% =============================================================================
\section{Monte Carlo Simulation}
\label{sec:monte-carlo}

Monte Carlo simulation estimates statistical properties of intractable
systems through repeated random sampling~\cite{metropolis1949monte,
robert2004monte}. In PATHCAST, MC generates synthetic process trajectories
by sampling from $P$ and $\hat{F}$, producing empirical distributions of
lead times and outcomes.

\subsection{Statistical Foundations}

\textbf{Law of Large Numbers.} For i.i.d.\ $X_1, X_2, \ldots$ with
$E[X] = \mu$:
\begin{equation}
\label{eq:lln}
\bar{X}_M = \frac{1}{M} \sum_{m=1}^{M} X_m \xrightarrow{a.s.} \mu
\quad \text{as } M \to \infty
\end{equation}

\textbf{Central Limit Theorem.} With $\text{Var}(X) = \sigma^2 < \infty$:
\begin{equation}
\label{eq:clt}
\frac{\bar{X}_M - \mu}{\sigma/\sqrt{M}} \xrightarrow{d} \mathcal{N}(0,1)
\end{equation}

The standard error scales as $O(1/\sqrt{M})$, dimensionality-independent.

\textbf{Estimator and confidence interval.} For a trajectory functional
$g$:
\begin{equation}
\label{eq:mc-estimator}
\hat{\mu}_g = \frac{1}{M}\sum_{m=1}^{M} g(\omega^{(m)})
\end{equation}
\begin{equation}
\label{eq:mc-ci}
\hat{\mu}_g \pm z_{\alpha/2} \cdot \frac{\hat{\sigma}_g}{\sqrt{M}}
\end{equation}

\subsection{Process-Aware Trajectory Sampling}

PATHCAST generates each trajectory by alternating: (i)~sojourn time draw
from $\hat{F}_s$, (ii)~next-state draw from $P[s, :]$, until an absorbing
state is reached. The full algorithm and convergence diagnostics appear
in Chapter~\ref{ch:method}, Algorithm~\ref{alg:mc-full}.

This formulation differs from the throughput-based MC approach of
Vacanti~\cite{vacanti2015}: throughput-based sampling treats the process
as a black box (aggregate output rate), whereas process-aware sampling
exposes state-specific delays, rework, and path frequencies. The
quantitative comparison is in Chapter~\ref{ch:method},
Table~\ref{tab:mc-comparison}.

\subsection{Reproducibility}

The Mersenne Twister PRNG (MT19937), default in NumPy and R, has period
$2^{19937}-1$ and passes the Diehard battery, far exceeding the
dimensionality of process simulations. Seed control makes MC outputs
exactly reproducible, satisfying the empirical software engineering
reproducibility requirement.


% =============================================================================
% 2.5 PROBABILISTIC FORECASTING METRICS
% =============================================================================
\section{Probabilistic Forecasting Metrics}
\label{sec:prob-metrics}

PATHCAST produces distributional forecasts; their evaluation requires
metrics that assess both point accuracy and distributional calibration.

\subsection{Continuous Ranked Probability Score}
\label{sec:crps}

\begin{definition}[CRPS]
\label{def:crps}
For a predicted distribution $\hat{F}$ and an observed value $y$:
\begin{equation}
\label{eq:crps}
\text{CRPS}(\hat{F}, y) = \int_{-\infty}^{\infty}
   [\hat{F}(t) - \mathbf{1}(t \geq y)]^2 \, dt
\end{equation}
\end{definition}

CRPS is a \emph{proper scoring rule}~\cite{gneiting2007strictly}, generalizes
the absolute error to distributions, and reduces to MAE when the forecast
degenerates to a point estimate. The CRPS--MAE relationship and
\emph{calibration error} derivations are in Appendix~H.

\subsection{Calibration}

A probabilistic forecast is calibrated if predicted quantiles match observed
frequencies. The Aggregate Calibration Error (ACE):
\begin{equation}
\label{eq:ace}
\text{ACE} = \frac{1}{|\mathcal{A}|} \sum_{\alpha \in \mathcal{A}}
    \left| \hat{p}_\alpha - \alpha \right|
\end{equation}
averages absolute deviation across decile checkpoints. ACE is reported
alongside CRPS because a forecast may be sharp but miscalibrated.

\subsection{Variant Entropy}
\label{sec:cap2-variant-entropy}

The Shannon entropy of the trace-variant distribution
$V = \{v_1, \ldots, v_k\}$ with frequencies $p(v_i)$ is
$H(V) = -\sum_{v} p(v) \log_2 p(v)$. The normalized form
$H_{\text{norm}}(V) = H(V) / \log_2 |V| \in [0,1]$ provides a
scale-independent diversity measure used in PATHCAST as a
process-complexity proxy. Formal definition and operational use are in
Section~\ref{sec:variant-entropy} (Chapter~\ref{ch:method}).


% =============================================================================
% 2.6 MACHINE LEARNING FOR SE
% =============================================================================
\section{Machine Learning for Software Engineering}
\label{sec:ml-se}

Machine learning has been extensively applied to defect prediction, effort
estimation, and code quality assessment~\cite{Wahono2015, pachouly2022}.
Tertiary studies~\cite{kotti2023} confirm that process-oriented prediction
tasks (throughput, workflow optimization) remain comparatively
underexplored.

\subsection{Supervised Learning for Process Prediction}

PATHCAST formulates outcome prediction as binary classification (on-time
vs.\ delayed) and lead time as regression, with feature vector
$\mathbf{x}_c \in \mathbb{R}^d$ and target $y_c$. The candidate algorithms
evaluated are Random Forests~\cite{breiman2001random},
XGBoost~\cite{chen2016xgboost}, logistic/linear regression baselines, and
LSTM networks~\cite{tax2017} for the sequence-based comparison. The
empirical evaluation protocol (temporal hold-out, hyperparameter search,
SHAP interpretation) is in Chapter~\ref{ch:methodology}.

\subsection{Process-Derived Feature Hypothesis}

The hypothesis underlying PATHCAST is that features derived from the
discovered process and the estimated Markov chain---conformance scores,
rework counts, expected remaining transitions, completion probability,
path entropy---add predictive value beyond traditional MSR features (code
metrics, issue attributes). The full feature catalogue and the comparison
protocol (PM-only vs.\ MSR-only vs.\ Combined) are in
Chapter~\ref{ch:method}, Table~\ref{tab:features}.


% =============================================================================
% 2.7 RELATED WORK AND GAP ANALYSIS
% =============================================================================
\section{Related Work and Gap Analysis}
\label{sec:related-work}

This section synthesizes the literature at the intersection of the five
domains, identifying the gap PATHCAST addresses. The systematic literature
review (Chapter~\ref{ch:slr}) provides the full evidence base; this section
distils the structural conclusion.

\subsection{Process Mining Applied to SDLC}

The principal contributions are summarized in Table~\ref{tab:pm-sd}: from
the seminal work of Cook and Wolf~\cite{cook1998}, through the MSR-PM
methodology of Poncin et al.~\cite{poncin2011process}, to recent industrial
applications by Caldeira and Cardoso~\cite{Caldeira2019AssessingSoftwareTeams}.
The consistent pattern is concentration at the descriptive and diagnostic
levels: discovery and conformance, rarely advancing to prediction or
forecasting.

\begin{table}[htbp]
\centering
\caption{Process mining applied to software development.}
\label{tab:pm-sd}
\small
\begin{tabular}{p{4.0cm}p{2.5cm}p{2.2cm}p{2.0cm}c}
\toprule
\textbf{Study} & \textbf{PM Type} & \textbf{Data Source} &
\textbf{SDLC Phase} & \textbf{Pred.} \\
\midrule
Cook \& Wolf (1998)            & Discovery       & Build logs       & Build/Test     & No \\
Rubin et al.\ (2007)            & Framework       & Conceptual       & All            & No \\
Poncin et al.\ (2011)           & Discovery       & VCS, Issues, ML  & Multi-repo     & No \\
Lemos et al.\ (2011)            & Disc., Conform. & Issue tracker    & Dev.\ lifecycle & No \\
Mannhardt et al.\ (2016)        & Conform.        & Multi-system     & Healthcare\textsuperscript{*} & No \\
Caldeira \& Cardoso (2019)     & Disc., Enhance. & Jira, Git        & Issue lifecycle & No \\
\bottomrule
\multicolumn{5}{l}{\textsuperscript{*}\scriptsize{Methodological reference for multi-perspective conformance.}}
\end{tabular}
\end{table}

\subsection{Stochastic Forecasting in SE}

Markov chains have been applied to software testing
(Whittaker and Thomason~\cite{WhittakerThomason1994}) and reliability
modeling~\cite{WhittakerRekabThomason2000}, but not parameterized by
process-mining-discovered structures.
Vacanti~\cite{vacanti2015} popularized throughput-based Monte Carlo for
agile forecasting, but this approach treats the process as a black box.
Rogge-Solti and Weske~\cite{RoggeSolti2015NonMarkovianSPN} integrated PM
with stochastic Petri nets for BPM remaining-time prediction, but did not
target SDLC.

\subsection{Machine Learning for SE Prediction}

Defect prediction~\cite{Zimmermann2005MiningVersionHistories, Wahono2015,
pachouly2022} relies overwhelmingly on artifact-level features.
PPM benchmarks~\cite{tax2017, marquezChamorro2018} operate on BPM event
logs, not software development. Camargo et al.~\cite{camargo2022learning}
combine deep learning with simulation but in BPM domains.

\subsection{Gap Synthesis}

Table~\ref{tab:gap} maps the intersection of techniques and domains.

\begin{table}[htbp]
\centering
\caption{Technique integration ($\checkmark$) across studies. PATHCAST is
the first proposal at the full intersection.}
\label{tab:gap}
\small
\begin{tabular}{p{3.8cm}ccccc}
\toprule
\textbf{Approach} & \textbf{PM} & \textbf{Markov} &
\textbf{MC Sim.} & \textbf{ML} & \textbf{SDLC} \\
\midrule
Cook \& Wolf (1998)             & $\checkmark$ &              &              &              & $\checkmark$ \\
Whittaker \& Thomason (1994)    &              & $\checkmark$ &              &              & $\checkmark$ \\
Vacanti (2015)                  &              &              & $\checkmark$ &              & $\checkmark$ \\
Zimmermann et al.\ (2005)        &              &              &              & $\checkmark$ & $\checkmark$ \\
Tax et al.\ (2017)               & $\checkmark$ &              &              & $\checkmark$ &              \\
Rogge-Solti \& Weske (2015)     & $\checkmark$ &              & $\checkmark$ &              &              \\
Rozinat et al.\ (2009)           & $\checkmark$ &              & $\checkmark$ &              &              \\
Camargo et al.\ (2019)          & $\checkmark$ &              & $\checkmark$ & $\checkmark$ &              \\
Marquez-Ch.\ et al.\ (2018)       & $\checkmark$ &              &              & $\checkmark$ &              \\
Caldeira \& Cardoso (2019)     & $\checkmark$ &              &              &              & $\checkmark$ \\
\midrule
\textbf{PATHCAST} & $\boldsymbol{\checkmark}$ & $\boldsymbol{\checkmark}$ &
$\boldsymbol{\checkmark}$ & $\boldsymbol{\checkmark}$ & $\boldsymbol{\checkmark}$ \\
\bottomrule
\end{tabular}
\end{table}

The gap decomposes into five deficits, each constituting a contribution
dimension: \textbf{(D1)~Descriptive Ceiling} (PM in SDLC stops at discovery);
\textbf{(D2)~Structural Blindness} (forecasting methods ignore process
structure); \textbf{(D3)~Stochastic Integration Absence} (Markov/MC not
parameterized by PM in SE); \textbf{(D4)~Feature Engineering Void} (ML
predictors do not incorporate PM-derived features); \textbf{(D5)~Pipeline
Fragmentation} (no formal end-to-end PM~$\to$~Markov~$\to$~MC~$\to$~ML
pipeline for SDLC). PATHCAST (Chapter~\ref{ch:method}) addresses all five
simultaneously through a four-stage pipeline with formal inter-stage
contracts.


% =============================================================================
% 2.8 CHAPTER SYNTHESIS
% =============================================================================
\section{Chapter Synthesis}
\label{sec:ch2-synthesis}

This chapter established the five theoretical pillars and provided the
formal vocabulary on which PATHCAST is built. \textbf{Process mining} provides
discovery, conformance, and enhancement; the discovered model is the
structural backbone. \textbf{MSR} provides the data sources and the
methodological principles for event-log construction.
\textbf{Markov chains} provide the formalism for state-transition dynamics,
with the absorbing-chain framework yielding closed-form expressions for
expected completion times and outcome probabilities. \textbf{Monte Carlo}
provides the computational engine for distributional forecasts. \textbf{ML}
provides predictive enhancement using process-derived features.

The integration of these five domains into a coherent pipeline with formal
inter-stage data contracts is the methodological contribution of
Chapter~\ref{ch:method}. The systematic literature review
(Chapter~\ref{ch:slr}) confirms that this integration has not been
previously proposed for software development. Detailed derivations and
algorithmic specifications are in Appendices~G--J:

\begin{itemize}
  \item Appendix~G: Markov chain derivations (proof of $N=(I-Q)^{-1}$,
  expected absorption time and variance, absorption probabilities,
  Bayesian smoothing, semi-Markov connection).
  \item Appendix~H: Probabilistic metric derivations (CRPS, CRPS--MAE
  relation, calibration error, worked examples).
  \item Appendix~I: Process mining algorithmic details (Inductive Miner
  internals, alignment computation, fitness/precision algorithms).
  \item Appendix~J: End-to-end illustrative walkthrough (event log $\to$
  process model $\to$ Markov chain $\to$ Monte Carlo simulation).
\end{itemize}
